\chapter{הפחתת סיכוניים וכיווניות עתידית}

\section{אסטרטגיות הפחתת סיכוניים}

על מנת להפחית סיכוניים בעת שימוש במודלים של שפה, חברות ועמותות חייבות לתכנן מסגרת שוטפת של בדיקה והשגחה. זה כוללת מספר אסטרטגיות:

\begin{enumerate}
    \item \textenglish{Red-teaming}: קבוצה של בודקים מנסים לחזות את הדרכים שבהן יכול מודל להיכשל, וכך לזהות חולשות.
    \item \textenglish{Continuous monitoring}: עקוב קבוע על התנהגות המודל בסביבה של ייצור כדי לזהות תופעות לא רצויות.
    \item \textenglish{Fallback mechanisms}: ניסיון לספק דרכים אלטרנטיביות או בדיקות מגניבה כאשר מודל לא בטוח.
    \item \textenglish{Fine-tuning and adaptation}: התאמה של המודל לסביבה ספציפית, באמצעות ייתרון או יישום של \textenglish{LoRA} (Low-Rank Adaptation).
    \item \textenglish{User feedback loops}: איסוף משוב מהמשתמשים הסופיים כדי לשפר את המודל וללמוד על סוגיות חדשות.
\end{enumerate}

\section{תקנון וחקיקה}

ככל שמודלים של שפה גדולים הופכים לחלק משכיח יותר של הציבור, ממשלות וגופים בינלאומיים מתחילים להתעניין בתקנון של תחום זה. כבר קיימות דוגמאות כגון \textenglish{EU AI Act}, אשר מעמיד כללים על מודלים של בינה מלאכותית בהקשר אירופאי.

תקנון יכול לכלול דרישות על:
\begin{itemize}
    \item תיעוד של נתוני ההדרכה ושיטות
    \item בדיקות של הטיות וביצועים
    \item הודעות שפופה למשתמשים על שימוש במודלים
    \item זכות למחוק מידע אישי מנתוני הדרכה
\end{itemize}

חברות כמו \textenglish{Anthropic}, \textenglish{OpenAI}, וגוגל משתתפות בדיונים אלה כדי להגדיר תקנים ושיטות בטיחות בתחום.

\section{כיווניות עתידית ומחקר פתוח}

התחום של בינה מלאכותית גדולה והשגחה עליה עדיין בשלבי ניצנים. מחקר עתידי יכול להתמקד בכמה תחומים:

\begin{enumerate}
    \item \textenglish{Interpretability}: הבנה של איך מודלים מקבלים החלטות ביותר עמוקה.
    \item \textenglish{Efficiency}: פתוח מודלים יותר קטנים וחסכוניים בחישוב שעדיין חזקים.
    \item \textenglish{Multilingual models}: משפרים קודים של מודלים שיכולים להתמודד עם שפות שונות, כולל שפות \textenglish{RTL} כמו עברית.
    \item \textenglish{Cross-modal learning}: שילוב של טקסט, תמונות, וסוגי מדיה אחרים.
    \item \textenglish{Human-in-the-loop systems}: מערכות שנותנות לאדם בקרה וביקורת על תוך כדי.
\end{enumerate}

מודלים כגון \textenglish{LLaMA}, \textenglish{Mixtral}, ו-\textenglish{Phi} מציגים טרנדים של דחיקה של גבולות היעילות. כלים כגון \textenglish{MCP} (פרוטוקול בקרת המודל) וטכניקות בקרת קשר מתקדמות יאפשרו יישומים ועסקים לנצל יכולות אלה בדרך בטוחה, חסכונית, ותמידית.

אנו בתחילת סוף מהפכה בתחום עיבוד השפה הטבעית, וה\textenglish{Transformer}-based מודלים הם בליבת מהפכה זו. התמודדות עם האתגרים של בטיחות, בקרה, ודיוק יהיה מרכזי לעתידות בוטחות ומשפחתיות של תחום זה.
